\subsection{Méthode}
\label{sec:Méthode}

\begin{frame}[c]{Méthode}
  \begin{columns}
    \column{.45\textwidth}
    \begin{itemize}
        \item Préciser le logiciel utilisé (ANSYS, Abaqus, ...)
        \item Préciser le type d'analyse (statique, modale, thermique, dynamique, ...)
        \item Décrire la démarche complète si elle dépasse la simple simulation (acquisition, identification, essais, \textbf{validation}, post-traitement, ...)
        \item Diagramme de flux pour les méthodes complexes (draw.io ou TikZ, si possible)
    \end{itemize}
    
    \column{.45\textwidth}
      \centering
      \begin{tikzpicture}[
        node distance=0.5cm,
        box/.style={draw,rectangle,minimum width=4cm,minimum height=0.8cm,align=center},
        arrow/.style={->,>=stealth}]

        % Nodes
        \node[box] (data) {
            Données\\
            $F,\ E,\ \nu$
        };
        \node[box, below=of data] (analytical) {
            Calcul analytique
        };
        \node[box, below=of analytical] (fem) {
            Simulation FEM
        };
        \node[box, below=of fem] (comparison) {
            Comparaison
        };
        \node[box, below=of comparison] (validation) {
            Validation
        };

        % Arrows
        \draw[arrow] (data) -- (analytical);
        \draw[arrow] (analytical) -- (fem);
        \draw[arrow] (fem) -- (comparison);
        \draw[arrow] (comparison) -- (validation);

      \end{tikzpicture}
  \end{columns}
\end{frame}

